\begin{statement}{B.36}{Proved}{The exact actual-source image on the abelian Fermat-two stage}
Retain the actual signed source $D$ of B.31 and the kernel
$D_0=\ker(D\to\operatorname{Aut}(\Delta))$ of B.33.  Put
\[
 E_2=C/K_2',\qquad W=K_2/K_2'\simeq\widehat{\mathbb Z}^{,5},
 \qquad H_{48}=C/K_2.
\]
Use the ordered Schreier basis $k_1,\ldots,k_5$ of B.24a and
write the group $W$ additively.  The primitive vector
\[
 w=(-1,1,0,1,0)
    =[yx^2y^{-1}]-[x^2]\in W                         \tag{B.36.1}
\]
has the following exact source interpretation.  There is a
continuous surjective homomorphism
\[
 \vartheta:D_0\longrightarrow\widehat{\mathbb Z},
 \qquad
 \vartheta(u)=-[k_1]\bigl(u(\rho)\rho^{-1}\bmod K_2'\bigr),
                                                        \tag{B.36.2}
\]
and every $u\in D_0$ acts on $E_2$ as
$\operatorname{Int}(\vartheta(u)w)$.  Here $[k_1]$ means
the first coordinate in $W$, and the inner automorphism uses
the actual element $\vartheta(u)w$ of this closed normal subgroup.
In particular
\[
 \ker(D\longrightarrow\operatorname{Aut}(E_2))
                 =\{u\in D_0:\vartheta(u)=0\}.          \tag{B.36.3}
\]

Let $\mathcal E_2$ be the nonsplit affine subgroup in B.35,
with its coordinates now denoted $(k,\mu)$.  Then the full image is
\[
 \begin{split}
 I_2:=\operatorname{im}(D\to\operatorname{Aut}(E_2))
   &\simeq\widehat{\mathbb Z}\rtimes\mathcal E_2\\
   &=\{(r,k,\mu):r,k\in\widehat{\mathbb Z},\
       \mu\in\widehat{\mathbb Z}^{\times},\
       k\bmod2=\eta_8(\mu)\},                           \tag{B.36.4}
 \end{split}
\]
where
\[
 (r,k,\mu)(s,l,\nu)
       =(r+\mu s,\ k+\mu l,\ \mu\nu).
                                                        \tag{B.36.5}
\]
The first factor is conjugation by $e_0^r$, and
$\mathcal E_2$ is exactly the image of the standard arithmetic
subgroup $G_{\mathbb Q}\to D$, measured by
$(\kappa_2,\chi)$.  Thus the kernel of its arithmetic action
is exactly $\{g:\chi(g)=1,\ \kappa_2(g)=0\}$.
For one fixed global path-orientation sign
$\varepsilon_{\rm path}\in\{1,-1\}$,
\[
 2\vartheta(g)=\varepsilon_{\rm path}\kappa_2(g)
 \quad\text{on }G_{\mathbb Q(\mu_\infty)}.              \tag{B.36.6}
\]
Changing the convention for a path cocycle changes this common
sign; neither the image nor any kernel above changes.
The projection $I_2\to I_1$ to the B.33 image forgets $k$,
and this extension does not split.

These are assertions about $C/K_2'$, whose kernel over $H_{48}$
is abelian.  They do not compute the full rank-five relative
pro-$\ell$ or prounipotent action, or the higher retained
normalizer images.
\end{statement}
\begin{proof}
\emph{The actual source fixes the six boundary classes.}
For $u\in D_0$, both $u(\rho)\rho^{-1}$ and
$u(\sigma)\sigma^{-1}$ lie in $F'\subset K_2$.  Thus $u$
induces the identity on $H_{48}$.  It fixes $x$ exactly.
Writing $h=u(\sigma)\sigma^{-1}\in K_2$ gives
$u(y)=hyh^{-1}$.  For the third peripheral element, first use
\[
 p=\rho^2x\rho^{-2}=x^{-1}y^{-1},\qquad
 z=(xy)^{-1}=xpx^{-1}.
\]
The element $u(\rho^2)\rho^{-2}$ also belongs to $F'$.
Consequently $u(p)$, and then $u(z)$, are conjugate to
$p,z$, respectively, by elements of $K_2$.
Since $u$ is the identity on $F/K_2$, it follows that it fixes
the six classes
\[
 [x^2],\ [yx^2y^{-1}],\ [y^2],\ [xy^2x^{-1}],\
 [z^2],\ [xz^2x^{-1}]                                \tag{B.36.7}
\]
in $W$.  These generate $W$ integrally, with their sum as the
sole relation.  This follows either from the six-punctured
genus-zero cover below, or from the integral vectors in B.24b:
their first five columns have determinant $\pm1$.
Hence $u|_W=1$.

\emph{The primitive integral line.}
For $h\in H_{48}$ choose any lift $s_h\in E_2$ and set
$b_h=u(s_h)s_h^{-1}\in W$.  The preceding paragraph makes
this independent of the lift and gives
\[
                  b_{hk}=b_h+A_hb_k,                    \tag{B.36.8}
\]
where $A_h$ is the actual component action on $W$.
The two generator matrices in the specified Schreier basis are
\[
 R_W=\begin{pmatrix}
 0&-1&-1&1&0\\0&0&-1&0&0\\1&0&0&0&-1\\
 0&1&0&-1&-1\\0&0&-1&1&0
 \end{pmatrix},\qquad
 S_W=\begin{pmatrix}
 0&0&1&0&0\\0&-1&0&1&1\\1&0&0&0&0\\
 0&0&0&0&1\\0&0&0&1&0
 \end{pmatrix}.                                         \tag{B.36.9}
\]
Let $P:W\to M$ be the matrix in (B.24a.1).  Direct integral
calculation gives
\[
 \begin{aligned}
 W^{e_0}&=\widehat{\mathbb Z}(1,0,0,0,0)
           \oplus\widehat{\mathbb Z}(0,1,0,1,0),\\
 W^{e_0}\cap\ker P&=\widehat{\mathbb Z}w,
       &W^{H_{48}}&=0,\\
 (1-R_W)w&=(1-S_W)w=(-1,1,1,1,-1),
       &A_yw&=-w.
 \end{aligned}                                         \tag{B.36.10}
\]
In the last equation $A_y=(R_WS_W)^2$.
In particular the $e_0$-fixed rank is two.  Its permutation
on the ordered six classes (B.36.7) is
$(0,1,4,5,3,2)$ in zero-based array notation, with cycle type
$1+1+4$.  Replacing this by the other degree-six $S_4$ action
would incorrectly discard the surviving line.

For each prime $\ell$, work temporarily in
$W_\ell\otimes\mathbb Q_\ell$.  Averaging (B.36.8) gives
\[
 t_\ell=\frac1{48}\sum_{h\in H_{48}}b_h,
 \qquad b_h=(1-A_h)t_\ell.                              \tag{B.36.11}
\]
There is no invariant vector, so $t_\ell$ is unique.
The fixed loop $e_0$ implies $t_\ell\in W^{e_0}_{\mathbb Q_\ell}$.
Since $u$ acts trivially on $\Delta$, $Pb_h=0$ for every $h$.
It follows that $Pt_\ell$ is invariant in $M_{\mathbb Q_\ell}$,
whose $H_{48}$ invariants vanish.  Thus $Pt_\ell=0$, and
(B.36.10) gives $t_\ell=\alpha_\ell w$.
The first coordinate of $(1-R_W)w$ is $-1$.  The integrality
of $b_\rho$ therefore forces $\alpha_\ell\in\mathbb Z_\ell$,
including at two and three.  Consequently the $t_\ell$ assemble
to the integral vector $\alpha w$, with
$\alpha=-[k_1]b_\rho\in\widehat{\mathbb Z}$.
Equation (B.36.8) is exactly the discrepancy of
$\operatorname{Int}(\alpha w)$.  This proves the asserted
formula for $u$, uniqueness, and the kernel (B.36.3).
It also proves that $\vartheta$ is a continuous homomorphism:
composition adds these unique inner translations.
Only the intermediate rational averaging used $1/48$;
the conclusion is integral because the displayed coboundary
vector is primitive.

\emph{A tangential Kummer calculation realizes the entire line.}
Use the arithmetic identification in B.35: $F$ is the
fundamental group of $U=\mathbb P^1-\{0,1,\infty\}$ based
at the tangent $z=0+1$, and $x,y$ are the inverses of its
two standard positive meridians.  The cover for $K_2$ is
\[
 u^2+v^2=1,\qquad z=u^2,
\]
with all points over $0,1,\infty$ removed.  Its rational
parameter and map to $U$ are
\[
 u=\frac{2T}{1+T^2},\qquad
 v=\frac{1-T^2}{1+T^2},\qquad
 z=\frac{4T^2}{(1+T^2)^2}.                             \tag{B.36.12}
\]
The six missing points are $0,\infty,1,-1,i,-i$.
The tangential lift with $u$-leading coefficient $1$ and
$v=1$ has $T$-leading coefficient $1/2$.  More explicitly,
write the local Puiseux parameter as $\epsilon$; then
$T=\epsilon/2$ gives $z=\epsilon^2+O(\epsilon^4)$.
This is the lift of the chosen tangent through the ramified
map, understood through its Puiseux fibre functor.

The deck element $y$ fixes $u$ and changes $v$ to $-v$.
In (B.36.12) it is $T\mapsto1/T$.
Thus the lift of the based loop $y$ is a path from the
tangent $T=0+1/2$ to the tangent at infinity with leading
coefficient $1/2$ in the coordinate $1/T$.
At the latter endpoint $T=2/\epsilon$ and again
$z=\epsilon^2+O(\epsilon^4)$.  Both tangential basepoints
are rational.  The arithmetic action from the braid source
is the action with this basepoint by B.35; the central
$\omega$ does not introduce a discrepancy inside $F$.

Let $G_*=G_{\mathbb Q(\mu_\infty)}$.  The action on
$W=H_1(U_{2,\overline{\mathbb Q}},\widehat{\mathbb Z})$
is trivial on $G_*$: all six punctures are cyclotomic and
their inertia generators are multiplied by $\chi=1$.
In particular changes of the lifted path do not change its
abelian cocycle on this group.
To compute that cocycle, use the five covers obtained by
adjoining an $m$th root of $T-a$, for
$a=0,1,-1,i,-i$, and all positive integers $m$.
The regularized endpoint values are the coefficients of
the lowest powers of $\epsilon$, rather than the values
of a function at a zero or a pole.  They give the table
\[
\begin{array}{c|ccccc}
 a&0&1&-1&i&-i\\\hline
 \text{initial value of }T-a&1/2&-1&1&-i&i\\
 \text{terminal value of }T-a&2&2&2&2&2\\
 \text{terminal/initial}&4&-2&2&2i&-2i
\end{array}                                             \tag{B.36.13}
\]
Indeed the first function has initial order one, the other
four have initial order zero, and all five have terminal
order minus one.

For clarity, the Kummer calculation here is integral at
every prime.  An $m$th root near either tangent is its
constant leading coefficient to the power $1/m$, times
a formal power of $\epsilon^{1/m}$, times a unit with
constant term one.  The latter unit has a unique root with
constant term one in characteristic zero.  Elements of
$G_*$ fix all roots of unity, all these formal parameters,
and every root of each root-of-unity factor in the final
row of (B.36.13).  The path's Galois discrepancy in this
cover is therefore the Kummer discrepancy of the ratio
of the two leading coefficients.  A change of path
multiplies its chosen transport by a root of unity, which
$G_*$ fixes.  Passing over all $m$ gives exactly, with
one common path-cocycle sign,
\[
                  (2,1,1,1,1)\kappa_2(g).              \tag{B.36.14}
\]
This is the tangential version of the elementary Kummer-cover
calculation in \citet[\S3.10 and Lemma 3.16]{FA_Wickelgren2012};
the direct leading-coefficient argument above also verifies
it at two, without using a prime-exclusion convention from
a nilpotent obstruction theorem.

If $b_a$ denotes the positive boundary class at $a$, the
pairing of $b_0-b_\infty$ with the five functions $T-a$
is $(2,1,1,1,1)$.  Since the original $x,y$ are negative
meridians, (B.36.1) identifies
$w=b_0-b_\infty$.  The cocycle of the lifted loop is the
class of $g(y)y^{-1}$, with the same fixed convention sign.
For $g\in G_*$, B.33 puts $g$ in $D_0$; the first part
of this proof computes the very same class as
\[
             (1-A_y)\vartheta(g)w=2\vartheta(g)w.
\]
The vector $w$ is primitive and the Kummer pairings are a
complete integral coordinate system on $W$.  Comparison
with (B.36.14) proves (B.36.6).
By (B.35.10), $\kappa_2(G_*)=2\widehat{\mathbb Z}$.
Multiplication by two is injective on $\widehat{\mathbb Z}$,
so $\vartheta(G_*)=\widehat{\mathbb Z}$.
This proves actual arithmetic realization of every integral
translation in the asserted line.

\emph{The full image and its nonsplit multiplier extension.}
Let $I_{\rm ar}$ be the image of $G_{\mathbb Q}$ on $E_2$.
Its action on $\Delta$ is $s_{\chi(g)}$ by B.33.
Consequently its action is trivial only if $\chi(g)=1$;
then (B.36.6) says that it is trivial exactly when
$\kappa_2(g)=0$.  Thus the action and the homomorphism
\[
 G_{\mathbb Q}\longrightarrow
       \widehat{\mathbb Z}\rtimes\widehat{\mathbb Z}^{\times},
       \qquad g\longmapsto(\kappa_2(g),\chi(g))
\]
have the same kernel.  Their images are isomorphic as
profinite groups.  The latter image is exactly $\mathcal E_2$
by (B.35.10)--(B.35.11) and cyclotomic surjectivity.
This argument identifies $I_{\rm ar}$ without asserting a
formula for the entire action from the multiplier alone.

For arbitrary $u\in D$, its first-stage image is uniquely
$\operatorname{Int}(e_0^r)s_\mu$.  Choose an arithmetic
element with multiplier $\mu$ and remove it together with
this inner loop factor.  The remaining action comes from
$D_0$ and, by the realization just proved, from an element
of $G_*$.  Thus $I_2$ is generated by the inner loop
subgroup and $I_{\rm ar}$.  They have trivial intersection:
comparison on $\Delta$ first forces $r=0$ and $\mu=1$,
after which the loop action is the identity.
The arithmetic subgroup conjugates
$\operatorname{Int}(e_0^r)$ to
$\operatorname{Int}(e_0^{\chi(g)r})$, by the defining exact
loop equation in $D$.  These facts prove (B.36.4)--(B.36.5),
including their profinite topologies by compactness.

A section of $I_2\to I_1$ would restrict over the subgroup
$(r,\mu)=(0,\mu)$ to a section of
$\mathcal E_2\to\widehat{\mathbb Z}^{\times}$.
B.35 proves that no such section exists, so this extension
is nonsplit as asserted.

The program \texttt{g\_fermat2\_source.py} uses the saved
integral extension of B.24a.  It verifies (B.36.9)--(B.36.10),
the unimodular boundary span, all $48^2=2304$ integral
cocycle identities for the surviving inner translation,
and its actual finite group reductions modulo $2,3,4$.
For each reduction it checks every element, every line
translation, and multiplication by both modular generators.
It separately verifies the rational parameterization, the
two pushed-forward tangents, and all five endpoint ratios
in (B.36.13).  The resulting exact certificate is
\texttt{g\_fermat2\_source.json}.  Arithmetic realization
uses the preceding path argument, not an assumption that
finite signed shadows lift to the source.
\end{proof}
